\section{Receiving Dionysius: Gregory and John of Damascus}
\label{sec:reception}

\subsection{Gregory: the orders within the joy of salvation}

Gregory the Great's thirty-fourth Gospel homily begins with the lost sheep and the lost coin. His account of the nine orders grows out of heaven's joy over the sinner's return. The angels and redeemed humanity belong to the one heavenly city; the ordering of that city leads the congregation toward repentance and charity. At \emph{Homilies on the Gospels} XXXIV.7, Gregory assembles the nine names from the prophetic books and Paul's epistles. At XXXIV.8--10 he explains their ministries. Angel designates a messenger's office; archangels announce greater things. Virtues work miracles, powers restrain hostile spirits, and principalities preside over good angels. Dominions stand above subordinate ranks; thrones serve divine judgment; cherubim excel in knowledge; seraphim burn with love.\footnote{Gregory the Great, \emph{Homiliae in Evangelia} XXXIV.3, 6--10, Latin text, Wikisource transcription of the complete homilies. The explanations here are paraphrases, not a received English translation.}


The parables give this teaching its center. In XXXIV.3, Gregory identifies the shepherd's lost sheep with humanity, borne home on the shoulders of the incarnate Lord. The faithful angels are the shepherd's friends because they keep his will, and his neighbors because they enjoy the vision of his glory. Christ asks them to rejoice with him: the recovered person's life is his joy. Angelic charity therefore follows the charity of Christ. The heavenly hosts welcome the restoration of those whose earthly weakness differs so greatly from their own.

The recovered coin develops the same mystery in XXXIV.6. Humanity bears the divine image; the lamp signifies the divine light made present in Christ's humanity; the house is stirred so that the lost image can be restored. Gregory associates the nine angelic orders and humanity with the ten coins of the parable. The governing claim is the gathering of rational creatures into the knowledge and love of God. This gives his treatment of the orders a pastoral urgency: the congregation is hearing about the city toward which it is being called, and about the charity with which that city already regards its return.

\subsubsection{The heavenly gifts in the life of the Church}

At XXXIV.11 Gregory asks what benefit the distinctions would have if they did not help his hearers grow. He then traces resemblances between angelic ministries and forms of Christian holiness. One believer faithfully teaches the little he has understood; another communicates deeper mysteries. Some receive gifts of miracles, others authority over hostile spirits, others responsibility for directing their brethren. Self-government gives an analogy to dominion, and a mind established in the fear of God can become a seat from which sound judgment serves others. These correspondences carry the account from heavenly order into the ordinary demands of teaching, discipline, and charity.

His treatment of cherubim is especially revealing. Fullness of knowledge finds its human likeness in fullness of love for God and neighbor, because charity fulfills the law. Understanding reaches its proper end in a rightly ordered life. His seraphic figures are those so inflamed by contemplation that their speech kindles love in others. Their ardor becomes illumination and purification: they help another see the good and repent of what obscures it. The preacher's account of the orders becomes an examination of how generously each hearer has received and communicated grace.

Gregory consequently warns against envy. The person with a smaller gift should seek growth and rejoice in another's greater gift. In XXXIV.14 the same principle explains the unity of the heavenly city: all the orders love, know, and receive God's rule, while a particular name marks the excellence especially abundant in one order. Through charity, each rejoices in the good of the others. The distribution of gifts enriches the common life. Aquinas receives this account in \ST{I q.~108 a.~5}; the names distinguish modes of excellence within a communion already perfected by love.\footnote{Gregory, \emph{Homiliae in Evangelia} XXXIV.3, 6, 11--14.}

His ascending list places virtues below powers and principalities above them. Dionysius places virtues above powers and principalities below them. Aquinas recognizes the exchange in \ST{I q.~108 a.~6}, relating Dionysius's sequence to Ephesians 1:21 and Gregory's to Colossians 1:16. He then argues that the difference diminishes when the functions attached to each name are compared.

Gregory also records a limit to his assent. Having reported Dionysius's account of the highest orders remaining within contemplation, he considers the seraph sent to Isaiah. In XXXIV.13 he declines to settle whether cherubim and seraphim act personally or through subordinate ministers when clear testimony does not decide it. He firmly affirms that angels send other angels, and that a mission does not deprive its bearer of inward contemplation. Aquinas adopts the more definite Dionysian account in \ST{I q.~112 a.~2}.

\subsection{John of Damascus: nature, grace, and an acknowledged unknown}

John's \emph{Exposition of the Orthodox Faith} II.3 gathers creation, freedom, illumination, location, ministry, and the nine orders into one compact account. Angels are created through the Word, perfected through the Spirit, and incapable of creating another essence. Their mobility and service express obedience to God. Their present stability in goodness is by grace and nearness to God, not the self-sufficiency of an uncreated nature. John expressly receives Dionysius's three triads.\footnote{John of Damascus, \emph{Exposition of the Orthodox Faith} II.3--4, trans. S.~D.~F.~Salmond, \emph{NPNF}, second series, vol.~IX (Edinburgh: T.~\& T.~Clark, 1898), CCEL transcription.}


\subsubsection{Intellectual life and obedient presence}

John describes the angels as secondary intelligible lights derived from the first light. Their knowledge is communicable without tongues or ears: one spirit makes its thought known to another through an intellectual act. The account joins immateriality with personal life. An angel has understanding and freedom, receives illumination, and wills the good it serves. Its obedience is an intelligent readiness. The imagery of fire expresses the warmth, keenness, and upward movement of this desire for God.

Their presence is likewise spiritual and finite. Walls and doors do not confine them, but only God is unlimited without qualification. An angel acts where it is assigned; its freedom from bodily extension does not give it divine omnipresence. John therefore speaks of their places according to the mode of an intellect's presence and operation. Aquinas develops a precise account of this relation in \ST{I q.~52 aa.~1--2}: angelic presence in a place is understood through the application of power, and a determinate created power does not occupy every place at once. The swift arrival of a heavenly messenger testifies to its obedience and ability, while the universality of providence belongs to God.

John gathers their ministries under divine praise and fulfillment of the divine will. Nations and regions are entrusted to their care; human affairs receive their help; divine mysteries are announced under forms suited to the recipient's sight. The same spirits who minister behold God according to their capacity, and John calls this vision their nourishment. Their service arises within a life satisfied by the divine good. Heaven's activity is consequently understood from its end: the spirits carry out God's will because their intellectual life is ordered toward him.\footnote{John of Damascus, \emph{Exposition of the Orthodox Faith} II.3.}

He calls the angel incorporeal, yet says that compared with God every creature is dense and material. Aquinas treats that relative expression in \ST{I q.~50 a.~1 ad~1}: even an incorporeal creature falls short of divine simplicity. A second difference is less readily reduced to vocabulary. John says that whether angels are equal in essence or differ is known to God, while recognizing their unequal illumination and rank. Aquinas argues positively that each angel differs specifically from every other (\ST{I q.~50 a.~4}). John's acknowledged unknown and Aquinas's philosophical determination have different force.

John favors the creation of angels before the visible world; Aquinas prefers their creation together with corporeal creatures as parts of one universe, while recognizing the contrary patristic opinion (\ST{I q.~61 a.~3}). In II.4 John makes the demons' wickedness a voluntary departure from an originally good nature. Their power remains subject to divine permission, and their predictions do not amount to omniscience. The same created nature that makes an angel excellent also sets the limits of a fallen angel's power. The evil will does not establish another god, another principle of being, or another providence.

\subsubsection{The fall and the boundary of hostile power}

John's account of evil in II.4 follows directly from this account of created goodness. The devil received a good nature and an honorable office; his rebellion was a departure from that gift. Darkness supplies the analogy because it is the absence of light. The rebel's will turns a good created power toward a disordered end, and other spirits follow him through their own choice. God remains the cause of their being and natural excellence, while the defect of their action belongs to the will that abandons its rule.

Their hostility has a real but bounded place in providence. John recalls the trials of Job and the Gospel account of the swine: the demons act under permission. Their knowledge also remains created. An angel announces the future when God reveals it; a demon may report an event observed at a distance or make an inference that proves false. John thereby distinguishes divine foreknowledge, communicated prophecy, and creaturely conjecture. He locates the practical consequence in fidelity: even a true prediction does not make a hostile spirit a trustworthy guide.

Temptation similarly does not compel consent. The attack can be received or resisted; its force does not make the tempter lord of the human will. Aquinas's account of demonic assault preserves the distinction between temptation and the inward movement of the will (\ST{I q.~114 aa.~1--3; q.~111 a.~2}). John closes by comparing the angels' fall with human death as a boundary beyond which repentance does not follow. Aquinas receives that comparison while supplying his own account of the firmness of an angelic choice (\ST{I q.~64 a.~2}). The conclusion belongs to a theology of judgment: created freedom is weighty because it is ordered to an enduring end, and divine grace perfects that freedom in the holy angels' stable adherence to God.\footnote{John of Damascus, \emph{Exposition of the Orthodox Faith} II.4.}
